\section{Attacchi a dizionario: principi e funzionamento}\label{sec:AttacchiDizionario}
% TODO: scrivere il contenuto del capitolo
Un attacco a dizionario è un attacco informatico in cui l'aggressore cerca di decifrare un codice cifrato o una password tramite una serie sistematica di tentativi di inserimento attingendo da uno o più dizionari di probabili chiavi.

Esso si basa su tre fasi:
\begin{enumerate}
    \item L'attaccante crea un dizionario, cioè una lista di possibili password. Spesso si inseriscono i nomi di animali domestici, di familiari o personaggi noti della cultura pop o di squadre sportive. Questo perché molte persone utilizzano parole di questo tipo per creare password che abbiano per loro un significato e siano facili da ricordare. Uno studio condotto da Google ha infatti rilevato che circa il 65\% delle persone utilizza le stesse password per più account e il 59\% utilizza nelle password dettagli personali facili da indovinare. L'elenco normalmente include diverse varianti, come varie combinazioni di parole oppure parole unite a caratteri speciali. 
    \item Il dizionario viene usato da un software automatizzato che tenta di violare gli account online.
    \item I dati sensibili a cui si ha accesso vengono sfruttati dall'attaccante per i propri scopi, per esempio per ottenere un guadagno economico o per perpetrare una frode.
\end{enumerate}

L'utilizzo di un dizionario differenzia questo metodo di attacco da quello a forza bruta. Quest'ultimo si basa sul tentare di scoprire una password provando tutte le combinazioni alfanumeriche possibili. Il vantaggio dato dal dizionario, quindi, consiste nel ridurre la quantità di combinazioni da provare, riducendo, di conseguenza, il tempo necessario, che in questi processi aumenta in maniera esponenziale all'aumentare della lunghezza della password stessa.
